\documentclass[12pt]{article}
\usepackage[utf8]{inputenc}
\usepackage[T1]{fontenc}
\usepackage{hyperref}
\title{Saraveca Language Resources for K-12 Curriculum Development}
\author{Compiled by Claude Sonnet 5.5 (claude-sonnet-5-5) for Andr\'as Kornai}
\date{}
\begin{document}
\maketitle

\section{Introduction}
Saraveca (ISO 639-3 \texttt{sar}; Glottocode sara1331) is an extinct Arawakan (Maipurean) language of Bolivia, spoken on the Rio Verde near the Brazilian border. Its closest relative is Paresi of Brazil. Sources disagree on when it died out; by 1962 most of the Sarave had shifted to Chiquitano. The language is known from about 400 words and phrases, plus a very short text and some small word lists. Any curriculum would be a heritage and revitalization programme built on this small record.

\section{Primary Educational Resources}
None found. There are no primers, textbooks or school programmes for Saraveca.

\section{Linguistic and Reference Materials}
\begin{itemize}
\item Cr\'equi-Montfort and Rivet (1913), the main source of about 400 words and phrases (part of their ``Linguistique bolivienne'' series); Glottolog reference record: \url{https://glottolog.org/resource/reference/id/152024}
\item Grambank entry for Saraveca: \url{https://grambank.clld.org/languages/sara1331}
\item Wikipedia overview, which lists the other records (Weddell 1851 text, d'Orbigny 1839 list of 23 words, unpublished Natterer and Fu\ss{} lists): \url{https://en.wikipedia.org/wiki/Saraveca_language}
\item Endangered Languages Project entry and bibliography: \url{https://www.endangeredlanguages.com/lang/2748?hl=pt-br}
\end{itemize}

\section{Religious and Community Resources}
None found.

\section{Digital and Audio Resources}
No audio exists for this language, as it has no known speakers recorded.

\section{Curriculum Development Context}
The documentation is limited to word lists and one very short text, so a K-12 curriculum cannot teach the language in a conversational sense. Realistic options are heritage-awareness units about the Sarave people and their history, and a small word-based component drawing on the 1913 list. Paresi, the closest living relative, might offer comparative material, but I found no resource for that. Descendants, who now mostly speak Chiquitano or Spanish, would need to be involved. I found no organisation representing them in the searches.

\end{document}